\begin{pastexam}[label={px:ccu-110-swa-08}]{中正資工所甲組 110 學年度 第 10 題}
（問答）

(3\%) If a dynamic programming problem satisfies the property of optimal substructure, justify if
a locally optimal solution is also globally optimal. A simple yes/no answer will not get any
point.
\end{pastexam}

\begin{pastsoln}
答案：不一定。 只有最佳子結構，局部最佳的選擇不保證得到整體最佳解；還要再加上貪心選擇性質才行。

\textbf{理由}。 最佳子結構說的是：整體的最佳解裡面，包含子問題的最佳解（資結書第~18.1.4~節）。
它的方向是「整體最佳 $\Rightarrow$ 拆出來的部分也最佳」，並沒有說「每一步挑當下看起來最好的，
最後就是整體最佳」。 動態規劃正是因為不知道哪一個選擇對，才把每一種選擇的子問題都解出來再比較。
「局部最佳就是整體最佳」需要的是貪心選擇性質：存在一個最佳解，它的第一個選擇就是當下最好的
那個選擇（資結書第~19.1.6~節）。 這是另一個性質，每個問題要另外證明。

\textbf{反例}：找零錢，面額 $1, 5, 12$，金額 $15$（資結書第~19.5~節）。
\begin{itemize}
\item 它有最佳子結構：令 $M[a]$ 是湊出 $a$ 的最少枚數，最佳解拿掉一枚面額 $c$ 的硬幣，剩下的一定
是湊出 $a - c$ 的最少枚數（不然換成更少的湊法，整體也更少），所以
$M[a] = 1 + \min_{c \le a} M[a - c]$，可以用動態規劃解，得 $M[15] = 3$（$5 + 5 + 5$）。
\item 每一步挑局部最好的（能拿的最大面額）：先拿 $12$，剩下 $3$ 只能用三個 $1$，共 $4$ 枚，不是
最佳。
\end{itemize}
另一種讀法是把「局部最佳解」當成子問題的最佳解，結論一樣：湊出 $12$ 的最佳解是一枚 $12$，但
$15$ 的最佳解 $5 + 5 + 5$ 根本不含它。 0/1 背包也是同樣的例子：有最佳子結構，但依單位價值貪心
不一定最佳（資結書第~19.2~節）。

反過來，像活動選擇、分數背包這些問題，除了最佳子結構還證得出貪心選擇性質，局部最佳的選擇
一路做下去才是整體最佳。
\end{pastsoln}
